\subsection{Entwickler-Onboarding}
\label{subsec:entwickler-onboarding}

Das Onboarding folgt dem Pfad Master klonen, \texttt{doctor}, \texttt{init},
Zielprojekt, erneut \texttt{doctor}, \texttt{install} und \texttt{run}.
Produktcode, Konfiguration, Tests und Commits gehören danach in das abgeleitete
Repository. Template-Wartung arbeitet dagegen im Master und prüft dort die
gesamte Baseline
\parencite{kleiveist2026governance-snapshot,kleiveist2026profiles}.
Diese Trennung verhindert, dass Produktentwicklung versehentlich die gemeinsame
Vorlage verändert.

Grundvoraussetzungen sind Git, Python ab 3.11 und Node.js ab 20. Desktopprofile
benötigen zusätzlich Rust Stable und plattformspezifische Tauri-Pakete;
Docker, PostgreSQL und PyGitIndex nur für die jeweiligen optionalen Aufgaben.
Optionale Werkzeuge blockieren damit nicht jedes Profil.

\path{AGENTS.md} stellt die Repository-Policy auch unterstützenden Coding
Agents bereit. Sie bezeichnet 900 Code-LOC als zulässige absolute Obergrenze,
fordert ab 600 LOC eine Prüfung und oberhalb von 750 LOC regelmäßig einen
Aufteilungsplan, untersagt die bloße Abschwächung von Regeln und verlangt dünne
Router, Adapter, Tauri-Kommandos und Composition Roots. Nach strukturellen
Codeänderungen soll vor der Übergabe das vollständige Quality Gate laufen;
anschließend sind relevante Tests auszuführen. Bei subsystemübergreifenden
Änderungen gilt die vollständige anwendbare Suite. Verbleibende Befunde sind
offenzulegen
\parencite{kleiveist2026agents}.

Die Agentendatei ist ein expliziter normativer Kontext, kein empirischer
Wirksamkeitsnachweis. Ihre Existenz belegt weder, dass ein konkretes Modell die
Regeln zuverlässig befolgt, noch dass agentengestützter Code dadurch
automatisch sauber ist. Technisch durchgesetzte Regeln, menschliches Review und
Modellverhalten bleiben getrennte Ebenen.

README und Fachdokumente trennen Einstieg, Tooling, Profile, Konfiguration, CI,
Container und Release. Das vereinheitlicht die öffentlichen Setup-Pfade,
belegt aber noch keine messbare Verkürzung des Onboardings
\parencite{kleiveist2026tooling,kleiveist2026configuration,kleiveist2026ci,kleiveist2026release}.
Die Dokumentationsstruktur ergänzt damit das technische Onboarding, ersetzt
aber keine empirische Messung seiner Dauer.
\beginnerinput{workspace/statement/500_tooling_workflow/570}
